\chapter{Συμπεράσματα}
\label{ch:conclusions}

Η εργασία αντιμετώπισε ένα πρακτικό ερώτημα της ελληνικής χονδρεμπορικής αγοράς: ποια
στρατηγική και ποιος τρόπος εκπαίδευσης παράγουν αξιόπιστη πρόβλεψη τιμής DAM και
ηλεκτρικού φορτίου κάτω από τους πραγματικούς περιορισμούς της Αγοράς Επόμενης Ημέρας,
όπου κάθε χαρακτηριστικό πρέπει να είναι γνωστό στο κλείσιμο της πύλης (D-1, 12:00 CET).
Το κεφάλαιο συγκεντρώνει τα ευρήματα, διατυπώνει τις συνεισφορές, οριοθετεί τους
περιορισμούς και προτείνει συγκεκριμένες επεκτάσεις.

% ─────────────────────────────────────────────────────────────────────────────
\section{Σύνοψη Κύριων Ευρημάτων}
\label{sec:conc_findings}
% ─────────────────────────────────────────────────────────────────────────────

Η εργασία αξιολόγησε τέσσερις στρατηγικές, δύο τρόπους εκπαίδευσης και πέντε
αλγορίθμους, για τιμή DAM και φορτίο. Ως προς τον σκοπό της, τα ευρήματα είναι τα εξής.

\paragraph{Στρατηγική, αλγόριθμος και τρόπος εκπαίδευσης.}
Το βασικό συμπέρασμα της εργασίας είναι ότι η επιλογή του πλαισίου εκπαίδευσης και της 
στρατηγικής πρόβλεψης επηρεάζει την ακρίβεια πολύ περισσότερο από την επιλογή του 
μεμονωμένου αλγορίθμου. Για την πρόβλεψη της τιμής, η μηνιαία επανεκπαίδευση μείωσε το μέσο απόλυτο 
σφάλμα κατά 27,7\% (από 14,478 σε 10,463 €/MWh για το LightGBM).
Στο ηλεκτρικό φορτίο, παρότι η επανεκπαίδευση προσφέρει οριακό μόνο όφελος ($\Delta\text{sMAPE} = -0,529\%$),
η Recursive πλησιάζει την Teacher-Forcing.
Οι αλγόριθμοι gradient boosting (LightGBM και XGBoost) προσφέρουν σταθερά την υψηλότερη ακρίβεια.
Ωστόσο, η πιο αξιόπιστη επιλογή είναι ο συνδυασμός τους (Ensemble-Best3 με βάση το $1/\text{MAE}$)
καθώς μειώνει το ρίσκο αστοχίας ενός μεμονωμένου αλγορίθμου.

\begin{table}[htbp]
\centering
\caption{Σύγκριση στρατηγικών Q1 2026 στον επιχειρησιακό ορίζοντα $H{=}24$ (ημερήσια
  δημοπρασία DAM).}
\label{tab:conc_summary}
\begin{tabular}{llrr}
\toprule
\textbf{Στρατηγική} & \textbf{Ορίζ.} & \textbf{Τιμή (€/MWh)} & \textbf{Φορτίο (MW)} \\
\midrule
Walk-Forward MR (TF Ensemble)$^{\dagger}$    & $H{=}24$ & \textbf{10,491} & \textbf{67,729} \\
\textbf{Walk-Forward Recursive}$^{\ddagger}$ & $H{=}24$ & 18,534          & 70,641 \\
Static TF (TF Ensemble)$^{\dagger}$          & $H{=}24$ & 14,478          & 96,751 \\
Static Recursive, best single                & $H{=}24$ & 20,094          & 102,961 \\
MIMO, best ensemble                           & $H{=}24$ & 20,896          & 217,088 \\
\midrule
Naive-24 (ρεαλιστική αναφορά)               & ---      & 21,638          & 343,130 \\
Naive-1 (αναφορά τύπου TF)                    & ---      & 15,433          & 267,151 \\
\bottomrule
\end{tabular}

\smallskip
{\footnotesize $\dagger$: oracle, μη ρεαλιστική (πραγματικές πρόσφατες υστερήσεις).\quad
$\ddagger$: επικρατούσα στρατηγική (\texttt{Rec-XGB-SS-DO-MR}).}
\end{table}

\paragraph{Δομικές Διαφορές και Συμπεριφορά των Σημάτων}
Παρά το κοινό πλαίσιο εκπαίδευσης και αξιολόγησης, τα δύο σήματα παρουσιάζουν εντελώς 
διαφορετική συμπεριφορά καθώς η τιμή είναι ένα έντονα στοχαστικό σήμα με συχνές αιχμές (spikes).
Αντίθετα, το ηλεκτρικό φορτίο είναι ένα ομαλό σήμα με ισχυρή περιοδικότητα και καθαρά ημερολογιακή δομή.  
Αυτή η δομική διαφορά καθορίζει σε μεγάλο βαθμό και την επίδοση των στρατηγικών πρόβλεψης. Στην περίπτωση 
του φορτίου, η αναδρομική Recursive καταφέρνει να πλησιάσει το θεωρητικό άνω φράγμα απόδοσης, ενώ στην τιμή απέχει αισθητά.  

\paragraph{Συμπεριφορά σε Περιόδους Έντονης Μεταβλητότητας}
Τον Φεβρουάριο του 2026 η τιμή έπεσε απότομα από τα $110 \pm 34$ €/MWh στα 57 €/MWh (-28\%) και η μεταβλητότητα εκτινάχθηκε στα $\pm$45 €/MWh. Τα στατικά μοντέλα βγήκαν εκτός του γνώριμου σε αυτά επίπεδο τιμών και έπεσε αισθητά η απόδοσή τους.
Το φορτίο την ίδια περίοδο έμεινε εντός των γνωστών ορίων ($5707 \pm 951$ MW η εκπαίδευση, $5567 \pm 839$ MW ο Φεβρουάριος), με μια απλή εποχική πτώση 7\%. Τα στατικά μοντέλα  απέδωσαν ικανοποιητικά και η επανεκπαίδευση προσέφερε οριακή βελτίωση ($\Delta\text{sMAPE} = -0,529\%$ στην Teacher-Forcing και μόλις $-0,15\%$ στην Direct).
Η Teacher-Forcing (που βασίζεται στις προηγούμενες τιμές) επηρεάστηκε τρομερά από τη βίαιη αλλαγή της τιμής, αλλά κέρδισε τα μέγιστα από την επανεκπαίδευση ($\Delta\text{sMAPE} = -4,496\%$). Αντίθετα, η Direct προσέγγιση λόγω των ανεξάρτητων μοντέλων ανά βήμα αποδείχθηκε πιο στιβαρή, με αποτέλεσμα μάλιστα στο σταθερό φορτίο η επανεκπαίδευση να της προσφέρει μηδαμινή βελτίωση ($-0,15\%$).

 \paragraph{Τεκμηρίωση σε τρία επίπεδα.}
Οι συγκρίσεις θεμελιώνονται ποιοτικά, ποσοτικά και στατιστικά. Στο ποιοτικό επίπεδο, η
ανάλυση της χρονικής δομής, της κατανομής και της μεροληψίας των σφαλμάτων αναδεικνύει τη
συστηματική υποεκτίμηση των αιχμών και τη διαφορετική συσσώρευση σφάλματος ανά στρατηγική.
Στο ποσοτικό, οι μετρικές MAE, RMSE και sMAPE δίνουν συμπληρωματική εικόνα (η Direct
συγκρατεί το RMSE έναντι της μεγαλύτερης διασποράς της MIMO στις αιχμές). Στο στατιστικό,
ο έλεγχος Diebold-Mariano με διόρθωση Harvey-Leybourne-Newbold επικυρώνει τις κατατάξεις:
υπεροχή της walk-forward TF στην τιμή ($p<0{,}001$), ενώ ισοδυναμία της αναδρομικής στο
φορτίο ($p=0{,}769$) και υστέρηση των MIMO/Direct έναντι κάθε μονοβηματικής στρατηγικής
($p<0{,}001$).

\paragraph{Επίδραση δεδομένων και σημαντικότητα χαρακτηριστικών.}
Η ποιότητα και η επεξεργασία των δεδομένων επηρεάζουν άμεσα την επίδοση. Οι πυκνές
ενδοημερήσιες υστερήσεις ωφελούν σταθερά τις πολυβηματικές στρατηγικές, όπου είναι
παρατηρημένες τιμές (MIMO: MLP $-22{,}3\%$, XGB $-7{,}9\%$), ενώ στην αναδρομική η επίδρασή
τους είναι περιορισμένη και εξαρτάται από τη χρήση προγραμματισμένης δειγματοληψίας.
Η ανάλυση σημαντικότητας λειτουργεί σε δύο επίπεδα. Ανά στρατηγική, καθορίζεται από τη διαθέσιμη
πληροφορία: το $y_{t-1}$ πέφτει από $88\%$ στο Teacher-Forcing σε $1$-$5\%$ στο MIMO, και
το κενό καλύπτουν μακρινές υστερήσεις, εξωγενείς παράγοντες στην τιμή και ημερολογιακά
χαρακτηριστικά στο φορτίο.
Ανά αλγόριθμο μεταβάλλεται ο τρόπος αξιοποίησης της ίδιας πληροφορίας. Τα μοντέλα ενίσχυσης 
κλίσης (gradient boosting) συγκεντρώνουν τη σημαντικότητα σε ελάχιστα χαρακτηριστικά,
με το LightGBM και το XGBoost να δίνουν στο \textit{Price\_lag1} $88{,}3\%$ και $63{,}3\%$ αντίστοιχα.
Αντίθετα, τα bagging μοντέλα κατανέμουν την πληροφορία ευρύτερα, με το Random Forest να δίνει
στο ίδιο χαρακτηριστικό μόλις $11,2\%$.

\paragraph{Προγραμματισμένη δειγματοληψία στο φορτίο.}
Στο φορτίο η προγραμματισμένη δειγματοληψία μειώνει σταθερά το σφάλμα: εκθέτοντας το μοντέλο
στις δικές του προβλέψεις κατά την εκπαίδευση, περιορίζει τη συσσώρευση σφάλματος και φέρνει
την αναδρομική στο επίπεδο του oracle. Στην τιμή το όφελος δεν είναι σταθερό και εμφανίζεται
κυρίως σε περιόδους μεταβολής καθεστώτος.

% ─────────────────────────────────────────────────────────────────────────────
\section{Συνεισφορές της Εργασίας}
\label{sec:conc_contributions}
% ─────────────────────────────────────────────────────────────────────────────

Η εργασία συνεισφέρει σε δύο επίπεδα. Σε μεθοδολογικό επίπεδο, ορίζει ένα ενιαίο πλαίσιο
χωρίς διαρροή δεδομένων που συγκρίνει τέσσερις στρατηγικές, δύο τρόπους εκπαίδευσης και
πέντε αλγορίθμους με κοινούς κανόνες, χρησιμοποιώντας μόνο πληροφορία διαθέσιμη στο κλείσιμο
της πύλης. Στηρίζει τις συγκρίσεις στον έλεγχο Diebold-Mariano με διόρθωση Holm-Bonferroni,
πρακτική σπάνια στην ελληνική βιβλιογραφία EPF/ELF. Μοντελοποιεί μαζί τιμή και φορτίο με κοινή
υποδομή, και προσφέρει ένα εργαλείο οπτικοποίησης που κάνει τη σύγκριση όλων των μοντέλων
προσιτή.

Σε εμπειρικό επίπεδο, μετρά το κέρδος της μηνιαίας επανεκπαίδευσης, δείχνει τα όρια της
θεωρίας συσσώρευσης σφάλματος για ομαλά σήματα όπως το φορτίο, περιγράφει πώς οι πυκνές
υστερήσεις δρουν διαφορετικά ανά στρατηγική, και χαρτογραφεί τη σημαντικότητα χαρακτηριστικών
σε δύο επίπεδα: ανά στρατηγική, όπου αλλάζει η διαθέσιμη πληροφορία, και ανά αλγόριθμο, όπου
αλλάζει ο τρόπος αξιοποίησής της.

% ─────────────────────────────────────────────────────────────────────────────
\section{Περιορισμοί}
\label{sec:conc_limitations}
% ─────────────────────────────────────────────────────────────────────────────

Η μελέτη έχει συγκεκριμένες αδυναμίες. Οι τιμές φυσικού αερίου για το διάστημα Νοεμβρίου 
2025 έως Φεβρουαρίου 2026 προέκυψαν με παρεμβολή από μεταγενέστερα σημεία, εισάγοντας 
εξάρτηση από το μέλλον στην τιμή αερίου της δοκιμαστικής περιόδου. Το SVR εκπαιδεύτηκε 
σε 5.000 από περίπου 80.000 δείγματα λόγω υπολογιστικού κόστους, οπότε τα αποτελέσματά του 
δεν είναι αντιπροσωπευτικά. Η θερμοκρασία, σημαντική για το φορτίο, μπαίνει μόνο έμμεσα 
μέσω ημερολογιακών χαρακτηριστικών και κυλιόμενων μέσων, ενώ η αξιολόγηση καλύπτει ένα μόνο 
τρίμηνο, αφήνοντας ανοιχτή τη γενίκευση σε άλλες εποχές. Η ωριαία συγχώνευση της παραγωγής με 
άθροισμα αντί μέσου όρου αλλάζει την κλίμακα της υπολειμματικής ζήτησης, καθώς το χαρακτηριστικό αυτό
δεν είναι από τα κύρια, η επίδραση κρίνεται μικρή αλλά χρειάζεται έλεγχο.
Παράλληλα, η ωριαία συγχώνευση των 15λεπτων δεδομένων εξομαλύνει την πραγματική ενδοημερήσια 
αστάθεια (intra-hour volatility) εντός της ώρας. Τέλος, το όφελος της προγραμματισμένης δειγματοληψίας
στο φορτίο και η επίδραση των πυκνών υστερήσεων στην αναδρομική δεν συνοδεύονται από πλήρες ablation, που 
μένει για μελλοντική δουλειά.

% ─────────────────────────────────────────────────────────────────────────────
\section{Μελλοντική Έρευνα}
\label{sec:conc_future}
% ─────────────────────────────────────────────────────────────────────────────

Οι προτεινόμενες επεκτάσεις πηγάζουν από τα ευρήματα και τους περιορισμούς.

Το σημαντικότερο περιθώριο βελτίωσης βρίσκεται στα δεδομένα. Η αντικατάσταση των συνθετικών
τιμών αερίου με πραγματικά ιστορικά αίρει την εξάρτηση από το μέλλον και επιτρέπει αξιόπιστη
επανεκτίμηση των μοντέλων τιμής. Η προσθήκη μετεωρολογικής πρόγνωσης (θερμοκρασία,
ηλιοφάνεια, ταχύτητα ανέμου) αντικαθιστά τα ημερολογιακά υποκατάστατα στα οποία στηρίζεται
σήμερα το φορτίο και αναμένεται να ενισχύσει τις πολυβηματικές στρατηγικές. Η ενοποίηση της παραγωγής σε μέσο όρο και ο επανυπολογισμός της υπολειμματικής ζήτησης
διορθώνει το ζήτημα κλίμακας, ενώ η εξέταση ιεραρχικών μοντέλων (hierarchical models) θα μπορούσε να αποτυπώσει τη 15λεπτη ενδοημερήσια αστάθεια.
Τέλος, η ακριβέστερη πρόβλεψη αιολικής και ηλιακής παραγωγής,
που η ανάλυση σημαντικότητας ανέδειξε ως κύριο εξωγενή παράγοντα της τιμής, αναμένεται να
ωφελήσει την πρόβλεψη τιμής.


Στο επίπεδο της μεθόδου, ο ρυθμός επανεκπαίδευσης μπορεί να βελτιστοποιηθεί: το επεκτεινόμενο
μηνιαίο παράθυρο μπορεί να συγκριθεί με συντομότερα κυλιόμενα παράθυρα και με online
ενημέρωση ημερήσιας συχνότητας, με στόχο ταχύτερη προσαρμογή στη μεταβολή κατανομής και
χαμηλότερο υπολογιστικό κόστος από την πλήρη επανεκπαίδευση.

Ξεχωριστό ζήτημα είναι η μοντελοποίηση των αιχμών. Όλες οι στρατηγικές υποεκτιμούν
συστηματικά τις ακραίες τιμές, οπότε μια ασύμμετρη συνάρτηση κόστους ή ένα ειδικό υπόδειγμα
για τις αιχμές θα αντιμετώπιζε στοχευμένα αυτή την αδυναμία.

Ευρύτερο ενδιαφέρον έχει η ζήτηση από κέντρα δεδομένων. Κατά τον Διεθνή Οργανισμό
Ενέργειας~\cite{iea2024_electricity}, η παγκόσμια κατανάλωσή τους προβλέπεται να
διπλασιαστεί προς το 2030 (περίπου 945~TWh, $\approx$3\% της παγκόσμιας ζήτησης),
με κύριο μοχλό την τεχνητή νοημοσύνη. Η
ρητή μοντελοποίηση αυτής της μεγάλης και σχετικά σταθερής νέας βάσης φορτίου, και της
επίδρασής της στην αιχμή και στη διαμόρφωση τιμής, είναι επίκαιρη επέκταση για την ελληνική
αγορά, όπου ήδη σχεδιάζονται αντίστοιχες εγκαταστάσεις.

% ─────────────────────────────────────────────────────────────────────────────
\section{Κλείσιμο}
\label{sec:conc_closing}
% ─────────────────────────────────────────────────────────────────────────────

Το πρακτικό συμπέρασμα είναι ότι η στρατηγική πρόβλεψης και ο τρόπος εκπαίδευσης έχουν
μεγαλύτερη βαρύτητα από την επιλογή αλγορίθμου: η τακτική ανανέωση των δεδομένων, και όχι
η πολυπλοκότητα του μοντέλου, καθορίζει την ακρίβεια στην ελληνική αγορά. Η μεθοδολογική
αυστηρότητα, η αποφυγή διαρροής δεδομένων, η πλήρης out-of-sample αξιολόγηση και ο
στατιστικός έλεγχος είναι η βάση κάθε αξιόπιστης σύγκρισης στη βιβλιογραφία EPF/ELF.

Οι συστάσεις διαφέρουν ανά χρήστη. Ο έμπορος ενέργειας, που χρειάζεται την τιμή στο κλείσιμο
της πύλης, μπορεί να υιοθετήσει αναδρομική με μηνιαία επανεκπαίδευση, με επιφύλαξη για την
υποεκτίμηση των αιχμών. Ο Διαχειριστής Συστήματος Μεταφοράς, που προβλέπει φορτίο, μπορεί να
βασιστεί στην αναδρομική με μηνιαία επανεκπαίδευση· οι στρατηγικές MIMO και Direct κρίνονται
ακατάλληλες. Όπου οι πραγματικές
πρόσφατες τιμές είναι διαθέσιμες κατά την πρόβλεψη, όπως στην ενδοημερήσια αγορά ή σε
λειτουργία πραγματικού χρόνου, το Teacher-Forcing παύει να είναι θεωρητικό όριο και
γίνεται και αυτό πρακτικά αξιοποιήσιμο.
